\section*{1273,15 К. Разрешаются ли две фазы?}
При \(N=96\) параметр порядка имеет два неглубоких максимума вблизи
\(\Delta\mu=3{,}46\)~эВ. Высота перемычки при равных высотах пиков
всего около \(0{,}64\,k_{\rm B}T\). На 288 узлах максимум и плечо в
этой области уже не дают двух равновесных бассейнов равной высоты.
В отличие от низкотемпературного результата, одно наличие плеча
в конечной ячейке не позволяет определить два предельных состава.

На 1400 узлах независимые траектории без зонтика, начатые из раствора
и упорядоченной структуры, при \(\Delta\mu=3{,}42\)~эВ дают
21,05 и 21,08~ат.\% Al, при 3,45~эВ --- 22,60 и 22,52~ат.\% Al.
В экспериментальном интервале 19,5--23,0~ат.\% Al максимум \(\chi_x\)
по перевзвешенной изотерме составляет \(\ChiMiddle\)~эВ\(^{-1}\) при
\(N=768\) и \(\ChiLarge\)~эВ\(^{-1}\) при \(N=1400\). В середине
интервала (21,25~ат.\% Al) средний глобальный порядок падает с 0,098
до 0,070 с увеличением ячейки от 768 до 1400 атомов.

Для проверки использована связь между откликом на химическое поле и
дисперсией состава, \(\chi_x=\partial\langle x\rangle/\partial\Delta\mu
=N\operatorname{var}(x)/(k_{\rm B}T)\). Если бы две фазы с разностью
составов 3,5~п.п. присутствовали с равными суммарными весами, один
межфазный вклад дал бы \(\chi_x=2{,}14\)~эВ\(^{-1}\) на 768 узлах
и 3,91 на 1400. Это оценка разрешимости при таком соотношении весов.
На 768 узлах флуктуации порядка максимальны около 27,62~ат.\% Al,
однако средние составы нижней и верхней четвертей распределения \(Q\)
там практически совпадают (27,59 и 27,63~ат.\%). На 1400 узлах при
3,56~эВ два старта дали 27,94 и 27,96~ат.\% Al, хотя средние \(Q\)
различались (0,090 и 0,263). Выше 30~ат.\% Al наибольшая ячейка
проверена прямыми точками около 30,4, 33,3 и 36,3~ат.\%. Перекрытие
между этими точками слабое. Узкий переход при
другом составе и характер потери дальнего порядка остаются открытыми.

В классическом C46 на 5120 узлах средние составы 51 окна покрывают
13,98--36,86~ат.\% Al. Поле возрастает между каждой парой соседних окон,
наименьший прирост составляет 6,53~мэВ. Во всех 384 блочных выборках
сохраняется монотонность. Внутренняя касательная с двумя контактами
не найдена ни в одной выборке. При перевзвешивании на квантовые
коэффициенты результат тот же. Эти данные дают независимую проверку
на уровне кластерной модели.
